\section{Manuscript audit, integration, and further research}
\label{ir:sec:audit}

\subsection{A concrete Clausen-parity correction}
The inspected chapter
\src{chapters/02-cyclotomic.tex} explicitly adopts
\[
 \Cl_{2m}(\theta)=\Im\Li_{2m}(e^{i\theta}),\qquad
 \Cl_{2m+1}(\theta)=\Re\Li_{2m+1}(e^{i\theta}).
\]
Here $\Cl$ denotes the manuscript's Clausen notation.
Later, under ``The fundamental constants, by conductor'', the sentence
``Odd Clausen values use the odd character sector'' is inconsistent
with this weight-index convention \cite{repo}. Odd-index Clausen values
are cosine-type real parts and use even characters, as the chapter's
own $\Cl_3(2\pi/5)$ formula with the even quadratic character already
illustrates.

A precise replacement is:
\begin{quote}
Imaginary, sine-type cyclotomic values use the odd character sector;
real, cosine-type values use the even sector. Under the Clausen
convention used here, even-index Clausen functions belong to the odd
sector and odd-index Clausen functions belong to the even sector.
\end{quote}
This is a terminology correction at the identified occurrence. It does
not invalidate the correctly stated Fourier formulas elsewhere in that
chapter. The integration directory supplies a guarded replacement script
and the exact old and new text; it does not modify the repository during
this delivery.

\subsection{An integral qualification, not a refutation of field results}
The manuscript's descriptions of Fourier-conjugate relations should
always retain their coefficient field. A Fourier transform is invertible
over a suitable cyclotomic field. It is not automatically a unimodular
change of integral lattices, since its inverse has a denominator.
Accordingly, agreement of field-valued ranks does not prove agreement of
integral Smith forms. This is a clarification of scope, not a claim
that the manuscript's explicitly field-valued identities are false.

The present calculations show why the qualification matters. Before
reflection, the chosen integral presentation has a unit maximal minor
and no torsion. After reflection, the quotient can contain a nonzero
$(\Z/2)^h$, although both characteristic-zero sign dimensions are
still $\varphi(q)/2$. In characteristic two, the rank jump is controlled
by the active ideal \eqref{ir:eq:universal-ideal}, not by a failure of
freeness of the unreflected distribution module.

It would therefore be inaccurate to summarize the result as
``distribution rank jumps at exceptional weights''. The correct
statement is that the \emph{reflection coinvariant} dimension can jump
in characteristic two. The original distribution module remains free
with the same raw-point basis at those weights.

\subsection{How to integrate the results without changing older claims}
A natural insertion is a section following the existing distribution
normal-form discussion, with cross-references from the surviving
research programme. The complete statement should include
Theorems~\ref{ir:thm:integral}, \ref{ir:thm:integer-reflection},
\ref{ir:thm:char2-koszul}, and \ref{ir:thm:smith-jets}. The
fixed-point reduction is short enough to retain rather than hide behind
an unexplained rank calculation. The analytic application can be
inserted beside the existing spectral-jet identities, with
\eqref{ir:eq:cert15} retained as a compact executable certificate.

The original characteristic-zero conductor basis remains useful and
correct within its stated scope. Its complex spectral weights should
not be reduced ``modulo two'' without first choosing a genuine integral
or modular coefficient model. The formal characteristic-two jets in
this article are power-series deformations of algebraic weights;
they are not reductions of real logarithms or complex spectral
parameters. Both uses of the word ``jet'' are explicit here, but their
coefficient categories are different.

The ordinary specialization is to be cited to the classical distribution
literature. The new integration should not relabel it as a first proof.
The existing mixed Gaussian $S_6$ conjecture retains its current status:
none of the depth-one module calculations in this article proves that
depth-two period identity. Similarly, no new numerical independence
claim follows from any formal free basis or nonzero module obstruction.

\subsection{Further research questions}
\paragraph{Integral equivariant structure beyond reflection.}
The proof treats one involution. For a subgroup of $(\Z/q\Z)^\times$
of higher order, determine the integral coinvariants and their primary
torsion. An appropriate orbit filtration would have to distinguish
several stabilizer types, rather than only fixed and moving points.
A first concrete target is a subgroup of prime order and a coefficient
characteristic equal to that prime. The raw basis and the finite
resolution give an explicit starting presentation.

\paragraph{A full equivariant normal form.}
The integral basis is deliberately not reflection-stable. Can one give
an explicit decomposition of the specialized lattice into trivial,
sign, and regular involution blocks, together with unimodular change
matrices of controlled size? The torsion calculation determines strong
constraints on such a decomposition, but the short fixed-complex proof
does not itself output those basis changes.

\paragraph{Multivariate deformation invariants.}
Theorem~\ref{ir:thm:char2-koszul} identifies every characteristic-two
base change through Koszul homology. Over a multivariate power-series
ring there need not be a Smith diagonal. Determine a useful replacement
in terms of Fitting ideals, minimal free resolutions, and intersection
multiplicities of the active ideal. The complete-intersection support
and the one-variable contact law provide tests for any proposed
multivariate formula. In particular, the universal cyclic homology
must not be naively specialized as if homology were flat.

\paragraph{Explicit transition to conductor-character coordinates.}
Construct the matrix between the raw-point basis and the manuscript's
conductor basis over cyclotomic integers localized at specified primes.
Determine its denominator ideal, not merely its determinant over a
splitting field. This would isolate exactly which arithmetic primes
obstruct character-coordinate replay and permit certified transport
between Fourier and raw-row certificates.

\paragraph{Shortest certificates and coefficient growth.}
The selected-row straightener gives one certificate, not necessarily
a shortest one. Determine coefficient-height and sparsity bounds as
functions of $q$, its prime factorization, and the target denominator.
At levels with many prime factors, compare different distinguished-lift
choices and pivot orders. A concrete optimization problem is to minimize
the number of raw rows for a fixed normal-form identity while keeping
integer polynomial coefficients.

\paragraph{Higher-depth distribution.}
At depth one, every root sum has a single target and the resolution
separates prime directions. Multiple polylogarithms introduce nested
index constraints, several colors, and shuffle/stuffle compatibility.
Define an integral higher-depth presentation before asking for its
rank or torsion. It would be especially useful to identify which
parts of the present unit-pivot argument survive and which introduce
new syzygies. This is a route toward stronger exact certificate
infrastructure, not a present proof of the $S_6$ conjecture.

\paragraph{Pole-aware identity generation.}
Use the degree bound \eqref{ir:eq:degree-bound} to generate families of
spectral derivatives at nonpositive integers and at the simple pole,
while keeping the endpoint contribution symbolic. Seek selections
whose real or imaginary parts remove most derivative coordinates.
The all-order family \eqref{ir:eq:jets15} is a small test case;
arithmetic reduction of its surviving period coordinates remains a
separate problem.

\paragraph{Certified numerics and formal proof.}
Replace the floating-point Mellin diagnostics by interval quadrature
with explicit endpoint and tail budgets. Independently, formalize
finite grids, the straightening termination order, and raw certificate
replay in Lean. The next algebraic layer is the split local complex
and finite-filtration proof; the periodic fixed-point reduction can
then be treated as a separate homological interface. These tasks would
convert precisely identified pieces of the present ordinary proofs
into machine-checked results.

\subsection{Conclusion}
The continuation resolves the integral and modular structure of the
specified depth-one weighted distribution presentation, while keeping
its analytic meaning and its limitations explicit. A basis can remain
free under every specialization even though a later symmetry quotient
acquires torsion. The fixed-point Koszul complex explains both the
integer parity rule and the full characteristic-two jet law. At the
analytic end, short polynomial row certificates give exact identities
at every spectral order, and their endpoint derivatives force a visible
Stieltjes correction. These are complementary uses of the same
integral presentation, rather than unrelated numerical discoveries.
